\documentclass[11pt]{article}
\usepackage[margin=1in]{geometry}
\usepackage[T1]{fontenc}
\usepackage{lmodern}
\usepackage{amsmath,amssymb}
\usepackage{booktabs,longtable,array}
\usepackage{hyperref}
\usepackage[numbers,sort&compress]{natbib}
\hypersetup{colorlinks=true,linkcolor=blue,citecolor=blue,urlcolor=blue}
\title{Forecasting Computer Science Breakthroughs in 2027--2031: An Evidence-Grounded Prospective Research Portfolio}
\author{Ada Mao \and GPT6-luna\\[3pt]\small Ada Mao: human research director; GPT6-luna: AI research collaborator}
\date{October 17, 2026\\[3pt]\small Version 1.0; evidence cutoff: October 17, 2026}

\begin{document}
\maketitle

\begin{center}
\begin{tabular}{ll}
Research type & Prospective Research Forecast / Research Portfolio / Methodological Study\\
Primary field & Computer Science\\
Secondary fields & Artificial Intelligence; Algorithms and Complexity; Scientific Forecasting\\
Research status & Qualitative, uncalibrated; not peer reviewed
\end{tabular}
\end{center}

\begin{quote}\textbf{Forecast status.} This paper reports a \texttt{QUALITATIVE\_UNCALIBRATED\_FORECAST}. It does not claim calibrated probabilities, validated Paper~3 performance on computer-science outcomes, Paper~4 planner superiority, independent human peer review, or that any forecast breakthrough has already occurred.\end{quote}

\begin{abstract}
Which computer-science research questions are most likely to see substantial, independently verifiable technical progress in 2027--2031? This paper presents an evidence-grounded prospective portfolio rather than a probability model. We screened the project's ten prospectively frozen qualitative research commitments, re-reviewed twelve earlier radar candidates, and assessed seven additional primary-source leads. We grouped dependent questions into research families, matched source claims to their stated models and endpoints, recorded counterevidence, and ranked opportunities separately for 2027--2028 and 2029--2031. The near-term portfolio prioritizes repository-scale code-and-proof synthesis, constrained large-language-model inference scheduling, intermediate-regime adaptive privacy filters, approximate fully dynamic maximum matching, and artifact-level multi-agent verification. The medium-term portfolio emphasizes the dynamic-matching/ordered Ruzsa--Szemer\'edi extremal chain, planar Three-Sets Cut-Uncut, planar discrete Fr\'echet distance under translation, two-stage integer-program optimization, and research agents that produce reproducible evidence. Several reported systems results make these areas concrete, but they are source-author results, not independent replications. The ten frozen commitments are distinct from the ten ranked portfolio entries, and the two rankings are not independent samples. Paper~3 has no eligible, independently adjudicated computer-science Target~A or Target~B evaluation; Paper~4 has no demonstrated action-selection advantage. We preserve an unresolved ORS deadline defect and an observation-window ambiguity in the three original events. The result is a transparent, falsifiable editorial portfolio, not evidence that AI can predict breakthroughs.
\end{abstract}

\noindent\textbf{Keywords:} computer science; research forecasting; prospective evaluation; evidence adjudication; formal verification; inference systems; dynamic algorithms.

\section{Introduction}
Research attention is easier to measure than scientific progress. Publication volume, citation growth, benchmark popularity, and topical similarity can rise without a theorem being proved, an algorithm improving against a matched baseline, or an empirical result surviving independent replication. Conversely, a technically important result may appear in a new formulation or in a venue that simple retrieval misses. Forecasting breakthroughs therefore requires an event specified before its outcome, a point-in-time evidence boundary, and a rule for distinguishing resolution from partial progress and missing evidence.

This paper addresses a narrower question than ``Can AI predict science?'': given the evidence available on 2026-10-17, which specific computer-science problem families merit the strongest qualitative expectation of independently verifiable progress by 2031, and what observations would change that ordering? We use ``progress'' in a technical sense and retain uncertainty. The portfolio is purposive and researcher-authored; it is not a representative sample of computer science.

The project's Paper~3 method studies question-level outcomes using frozen, point-in-time features and temporal validation. Its astronomy results cannot be transferred to computer science without eligible CS outcomes. Paper~4 organizes competing hypotheses, evidence gaps, and candidate research actions; its utility estimates are not observed scientific value. We use these methods as design constraints and vocabulary, not as validated forecasting or planning engines in this study.

\section{Background and related work}
Science-of-science research has forecast citation impact and bibliometric ``breakthrough'' proxies, including work using early citation-network structure \citep{min2021breakthrough}. PreScience provides a large, time-snapshotted scientific forecasting benchmark for paper- and topic-level outcomes \citep{ajith2026prescience}. Neither establishes the exact task here: forecasting whether a frozen, theorem-level computer-science question will be resolved or substantially advanced under predeclared assumptions. This distinction does not make this portfolio the first scientific forecasting benchmark.

The project's question-discovery and astronomy historical-backtesting work is closer in unit and workflow: freeze a question and evidence cutoff, search later literature, and adjudicate outcomes \citep{mao2026questiondiscovery,mao2026backtesting}. Its prospective astronomy instance is same-project prior work, not independent corroboration. It also documents a serious challenge: outcome categories may be difficult for careful reviewers to apply consistently. ResearchBench evaluates inspiration retrieval, hypothesis composition, and ranking \citep{liu2026researchbench}; PaperBench evaluates replication of published AI research \citep{starace2025paperbench}. These are relevant but distinct from prospective theorem-resolution forecasting.

\subsection{Paper~3 and Paper~4}
Paper~3's baseline uses question-by-cutoff observations, features limited to information available at each historical cutoff, and temporal rather than random validation. Future attention, direct resolution, and premise revision are distinct outcome dimensions. Its quantitative astronomy attention results do not imply CS predictive skill. Here we preserve two targets separately: \textbf{Target~A}, direct proof or refutation of the frozen proposition, and \textbf{Target~B}, substantial independently verifiable progress even if the proposition remains open. Target~B may include a verified theorem or bound, matched-baseline algorithmic improvement, independent replication, independently checked counterexample, justified premise revision, or substantive narrowing. Attention is not a success label for either target.

Paper~4 represents a question through competing hypotheses, a provenance-linked evidence graph, missing or contradictory support, candidate actions, and declared cost and information-gain assumptions. That structure is useful for documenting what evidence would change a ranking. It does not show that the highest-ranked action is optimal or that the planner outperforms simple baselines. No information-gain probabilities or efficacy results are inferred here.

\section{Research design}
\subsection{Evidence boundary and selection}
The evidence cutoff is \textbf{2026-10-17}. We reviewed the final portfolio dossier, the ten-record prospective cohort and its manifest, earlier CS radar and source-screening material, Paper~3 and Paper~4 methodology records, and novelty/readiness audits. We screened primary-source leads in formal verification, inference systems, adaptive privacy, dynamic algorithms, approximation, geometry, integer programming, and scientific-agent evaluation. Candidate claims were checked against the source's actual graph/problem class, parameterization, baseline, metric, and assumptions. Where a reported benchmark result is not independently reproduced in this project, it is explicitly attributed to its authors.

We grouped related questions before ranking. In particular, approximate dynamic maximum matching and the ordered Ruzsa--Szemer\'edi (ORS) extremal parameter form one dependent chain, not two independent confirmations. We did not treat absence of a retrieved paper as evidence that a question is neglected or unresolved.

\subsection{Ranking method}
The portfolio contains five entries for 2027--2028 and five for 2029--2031. Each entry had to have a source-backed baseline, a technically interpretable progress endpoint, an identifiable obstacle, counterevidence, and a plausible independent verification route. We then made a qualitative ordering using evidence and endpoint clarity, plausible technical path, horizon tractability, potential significance, and counterevidence. We used no numerical weights, probabilities, attention counts, or citation counts. The order is a decision-priority judgment by the research authors, not an estimate of success probability.

The final portfolio is separate from the ten frozen commitments described in Section~\ref{sec:registry}. Seven commitments were added to the original three in a prior registration step. The final 5+5 rankings neither create ten new events nor replace the frozen texts. No new event passed the threshold and independence gate in the final portfolio round.

\section{What counts as a breakthrough or progress?}
We use ``breakthrough'' in the title as shorthand for substantial technical progress on a specified research question; it is not a formal label assigned to the listed opportunities. A qualifying progress observation must be independently verifiable and may be: (i) a theorem or bound advancing the exact stated model; (ii) an algorithmic improvement against a matched, explicitly stated baseline; (iii) an empirical improvement independently reproduced under a comparable protocol; (iv) an independently checked counterexample to a stated claim; (v) a justified refutation or revision of a key premise; or (vi) a material narrowing of the open question that rules out a meaningful class of approaches.

These categories can coexist and should not be collapsed without a separate protocol. Direct proof or refutation of the frozen proposition is Target~A. Substantial progress short of full resolution is Target~B. Attention, citation or publication activity is tracked separately; practical computational impact is also separate. Missing or inaccessible evidence is unadjudicated, not a negative outcome. An algorithmic proof establishes correctness relative to its formal specification, not that the specification captures intended behavior.

\section{Near-term portfolio: 2027--2028}
\subsection{Repository-scale code-and-proof synthesis}
\textbf{Question.} Can automated systems produce buildable implementations and machine-checked proofs across multi-module repositories with cross-module invariants, rather than solve isolated proof obligations?

\textbf{Baseline and evidence.} Vero reports 43 repository instances; its authors report 27 complete code-and-proof solves for the strongest evaluated configuration and ten instances resistant to all tested configurations \citep{ye2026vero}. VeriSoftBench reports difficulty with deep repository dependencies \citep{xin2026verisoft}. CoCo-Prover reports improved success-cost results, including up to 26/43 Vero solves, in its own evaluation \citep{yao2026cocoprover}. These are author-reported benchmark results; this project has not reproduced them.

\textbf{Expected progress and obstacle.} The qualitative expectation is further benchmark-specific improvement in solve coverage and cost-quality trade-offs, not reliable autonomous software engineering in general. Cross-module invariants, reusable proof structure, complete builds, benchmark curation, shifting models and compute budgets, and correspondence between formal specification and intended behavior are major obstacles. Verify on frozen repositories, hidden held-out tasks, trusted proof-kernel checks, clean builds, independent specification audits, and equal compute/cost budgets.

\subsection{LLM inference scheduling under KV-cache pressure}
\textbf{Question.} Can a serving policy jointly schedule requests, speculative work, and KV-cache residency under changing arrivals to improve goodput or tail latency at fixed output quality, memory, and service constraints?

\textbf{Evidence and obstacle.} Online scheduling results depend on arrival assumptions; one fully adversarial model rules out a constant competitive ratio \citep{jaillet2025kv}. Spexis reports up to a 34\% request-rate increase on evaluated GPU/workload configurations \citep{jung2026spexis}; it is not independently replicated here. Gains may disappear when memory, model quality, workload, preemption, energy, and p95/p99 latency are matched. Replay fixed arrival traces on matched hardware and model; report goodput, tail latency, energy, eviction/recomputation, code, and failures. Separate empirical systems gain from an online theorem.

\subsection{Adaptive privacy filters at intermediate sampling rates}
\textbf{Question.} Can valid, quantitatively useful filters be proved for fully adaptive compositions of Poisson-subsampled Gaussian mechanisms in intermediate sampling-rate regimes not covered by current approximate constructions?

Tran et al. show that a natural fully adaptive $f$-DP filter can be invalid, state validity conditions, and construct approximate GDP filters in low-rate $q<0.2$ and high-rate $q>0.8$ regimes under their assumptions \citep{tran2026fdp}. This makes the middle-regime question precise but does not establish that a useful extension exists. Full adaptivity, approximation error, and numerical soundness are difficult; existing RDP or privacy-loss-distribution accountants may be stronger in some settings. The frozen event's comparator remains flagged as underspecified. Verification requires a proof with explicit assumptions and error bounds, and comparison against named accountants under identical composition models.

\subsection{Approximate fully dynamic maximum matching}
Can a general-graph fully dynamic data structure maintain an explicit $(1-\varepsilon)$-approximate maximum matching in $O(n^{1-\delta})$ amortized update time for fixed $\varepsilon,\delta>0$, under the frozen adversary and output assumptions? FOCS~2024 and SODA~2025 work connects this update time to ordered-RS density \citep{behnezhad2024dynamic,assadi2025ors}. Bernstein et al. later show sublinear update time for \emph{maximal} matching \citep{bernstein2025maximal}; maximal matching is not approximate maximum matching. A better extremal bound helps only if the reduction preserves the desired exponent and all approximation, adversary, amortization, and explicit-output guarantees.

\subsection{Artifact-level verification of multi-agent work}
Can a verifier predict independently checked final-artifact correctness and improve end-to-end task success after accounting for cost, false positives and negatives, and context? MAS-ProVe reports that process verification does not consistently improve multi-agent performance and may have high variance \citep{venkataramani2026masprove}. EIG and ResearchBench assess ideation; PaperBench tests replication. These are not a common independent-discovery outcome. Freeze verifier outputs before independent artifact adjudication and report error rates, agreement, compute, and paired task success. A judge score is not itself a correctness label.

\section{Medium-term portfolio: 2029--2031}
\subsection{Dynamic matching and the ORS extremal bottleneck}
The extremal question is whether ordered Ruzsa--Szemer\'edi at fixed linear matching density has polynomial growth strong enough to imply a fixed-exponent sublinear update bound for approximate maximum matching. This is one dependent research family, not two observations. An ORS improvement must transfer through the exact reduction and its quantifiers. The frozen ORS event has an unresolved date conflict (Section~\ref{sec:registry}).

\subsection{Planar Three-Sets Cut-Uncut}
Can unweighted planar Three-Sets Cut-Uncut with two singleton terminal sets and one two-vertex terminal set that must remain connected be solved exactly in polynomial time? Bentert et al. give a randomized fixed-parameter algorithm for the planar two-set version and explicitly leave the $(1,1,2)$ three-set case open \citep{bentert2024twosets}. Ordinary planar Multiway Cut does not enforce the required connectivity. Verification must compare terminal pattern, graph class, costs, connectivity, exactness, and infeasibility handling.

\subsection{Planar discrete Fr\'echet under translation}
The frozen question targets exact discrete Fr\'echet distance under translation between planar polygonal curves in $O(n^4\operatorname{polylog}n)$ time under a comparable model. The audited source gives an approximately $\widetilde O(n^{14/3})$ algorithm and SETH-based $n^{4-o(1)}$ lower bound for the specified problem \citep{bringmann2019frechet}. The remaining exponent room is narrow and conditional; computational model, arithmetic, exact-value versus decision, and bit complexity must be matched.

\subsection{Two-stage integer programming}
Can the FPT feasibility algorithm for two-stage block-structured integer programs extend to exact linear optimization when local coefficient parameters remain bounded but global-block coefficients are unrestricted? Cslovjecsek et al. establish neighboring feasibility results and identify optimization as a distinct next question \citep{cslovjecsek2025blockip}. A later FPT theorem parameterized by all coefficients does not settle the unbounded-global-coefficient regime. The endpoint must distinguish exact optimization, infeasibility, and unboundedness under the same encoding and parameters.

\subsection{Scientific agents producing reproducible evidence}
Can agents freeze a falsifiable question, execute a discriminating study, preserve raw artifacts, and obtain independent reproduction more reliably than simple baselines? Current tools separately measure ideation, replication, or process verification; they do not establish a common outcome for original, independently confirmed science. This is the least mature entry. A useful advance requires a reliable adjudication taxonomy, agreement and missingness reporting, raw-artifact reproduction, and simple baselines. No new forecast event is frozen for this family.

\section{Existing prospective forecast registry}\label{sec:registry}
The registry contains \textbf{ten frozen qualitative commitments}: three original records and seven later registrations. They are commitments within a research program, not ten newly generated portfolio rankings and not ten independent observations. The complete immutable text, source URLs, event definitions, source cutoffs, issue timestamps, and hashes are at \url{https://github.com/nonameisready/scientific-discovery-computer-science/blob/db6aaa4b7f5fa364ced3042bf55d8df6b8eef1ce/data/forecast_engine/prospective_cohort_v2.json}; the linked activation record preserves the original three.

\small
\begin{longtable}{p{0.28\linewidth}p{0.16\linewidth}p{0.14\linewidth}p{0.32\linewidth}}
\toprule Frozen event & Issue UTC & Deadline & Compact endpoint / primary source \\
\midrule\endfirsthead
\toprule Frozen event & Issue UTC & Deadline & Compact endpoint / primary source \\
\midrule\endhead
\texttt{PLANAR-THREE-SETS-CUT-UNCUT-112} & 2026-10-17 02:47:37 & 2030-12-31 & Exact planar $(1,1,2)$ Cut-Uncut algorithm; Bentert et al. (ICALP 2024). \\
\texttt{DISCRETE-FRECHET-TRANSLATION-PLANAR-N4} & 2026-10-17 02:47:37 & 2031-10-16 & Exact planar discrete Fr\'echet under translation; Bringmann et al. \\
\texttt{K-HOP-DOMINATION-LOW-EXPANSION-NO-GIRTH} & 2026-10-17 02:47:37 & 2031-10-16 & Remove high-girth restriction at same stated bound; Lenzen and Wenning. \\
\texttt{MAXBISECTION-GW-APPROX-001} & 2026-10-17 03:15:09 & 2031-12-31 & MAX BISECTION at the Goemans--Williamson MAX CUT ratio; Brakensiek et al. (SODA 2026). \\
\texttt{DYNAMIC-MATCHING-POLY-EXPONENT-001} & 2026-10-17 03:15:09 & 2031-12-31 & Explicit approximate maximum matching with fixed-exponent sublinear dynamic updates; Behnezhad and Ghafari. \\
\texttt{ORS-POLYNOMIAL-GROWTH-001} & 2026-10-17 03:15:09 & \textbf{2030 structured / 2031 event text} & Polynomial ordered-RS growth at linear density; Assadi et al. \\
\texttt{SUBSET-SUM-EXPONENT-IMPROVEMENT-001} & 2026-10-17 03:15:09 & 2031-12-31 & Fixed-base exponent improvement for worst-case exact Subset Sum; Chen et al. (ESA 2025). \\
\texttt{TWO-STAGE-IP-OPT-UNBOUNDED-GLOBAL-COEFF-001} & 2026-10-17 03:15:09 & 2030-12-31 & FPT exact optimization with unrestricted global coefficients; Cslovjecsek et al. \\
\texttt{SHALLOW-QC-IDENTITY-PTAS-001} & 2026-10-17 03:15:09 & 2031-12-31 & PTAS for diamond distance of constant-depth local circuit from identity; Bravyi et al. \\
\texttt{ADAPTIVE-FDP-GDP-FILTER-MID-Q-001} & 2026-10-17 03:15:09 & 2030-12-31 & Valid approximate fully adaptive GDP filter at intermediate rate; Tran et al. \\
\bottomrule
\end{longtable}
\normalsize

The original three records share the issue timestamp 2026-10-17 02:47:37 UTC, but their individual \texttt{observation\_start} fields are null. Their issue date precedes the global 2027-01-01 horizon start, so whether results in the intervening period count is ambiguous. We preserve rather than retroactively repair this. The later seven records were issued at 03:15:09 UTC and record source-only cutoffs at issuance. Final adjudication requires two independent human reviewers; AI review is not independent human review. Source hashes cover citation/metadata records and registry snapshots, not full-text bytes.

The ORS record has an unresolved integrity defect: the structured deadline is 2030-12-31, while frozen event text states 2031-12-31. We do not select a preferred deadline or rewrite the registry. If the conflict affects adjudication, the outcome must remain unresolved pending an authorized provenance-based correction.

\section{Ranking robustness}
We compared three qualitative priority scenarios while holding candidate families fixed. Proof-tractability priority raises adaptive privacy, planar cut-uncut, and two-stage IP. Computational-impact priority raises inference scheduling and repository-scale code/proof. Source maturity and independent verifiability raises theorem endpoints and lowers process verification and autonomous science, whose outcome labels are least settled. Leading family membership is more stable than exact order; medium-term ranking is especially sensitive. These are ordinal comparisons, not statistical estimates, calibration intervals, or probabilities.

\section{Scientific validation framework}
Future monitoring should preserve exact questions, source cutoffs, issue timestamps, family identifiers, and predeclared outcomes. Theorem matching requires reviewers to compare the object, assumptions, graph/problem class, model, parameters, quantifiers, and endpoint. Empirical verification requires versioned artifacts, matched baselines, comparable resources, uncertainty, and reproduction. Reviewers should record disagreement; search failure remains missing evidence.

For Target~A, evaluate direct proof/refutation. For Target~B, separately adjudicate verified progress without collapsing it into resolution. Attention and impact remain separate. A future Paper~3 comparison requires sufficient eligible outcomes, point-in-time features, family-grouped chronology, an untouched final period, and baselines such as publication volume, field momentum, and source-only text. No CS score is valid before those gates pass. Paper~4 should timestamp competing hypotheses, provenance, actions, costs, and information-gain assumptions before action; afterward separate feasibility from evidence yield. The present portfolio measures neither predictor performance nor planner efficacy.

\section{Limitations}
This is one purposive qualitative synthesis, not a representative cohort or calibrated model. The ten commitments are small and heterogeneous, and most mature only in 2030--2031. No outcome labels have matured; independent human event adjudication and independent panel review are incomplete. Candidate families are dependent; publication, selection, and reporting bias are plausible. Vero, CoCo-Prover, and Spexis measurements are author-reported and not independently reproduced here. Some empirical topics need more precise common endpoints. Source hashes cover metadata, not full text. The original three have an observation-window ambiguity and ORS has a deadline mismatch. No calibrated forecasting accuracy, Paper~3 CS predictive performance, or Paper~4 action-selection superiority has been demonstrated. This paper is not peer reviewed.

\section{Conclusion}
The evidence supports a ranked set of questions that merit attention as prospective research opportunities, not claims of predictive success. Repository-scale code-and-proof synthesis and constrained inference scheduling have the strongest qualitative near-term case because they have visible technical progress and objective artifact-level tests, yet reported gains remain configuration-sensitive. Adaptive privacy filters offer a precise theory question. Approximate maximum matching is linked to the unresolved ORS bottleneck; maximal matching progress is not a substitute. The medium-term portfolio includes theorem-level planar, geometric, and integer-programming questions, while reproducible scientific agents remain uncertain because outcome adjudication is itself unresolved.

The ten frozen commitments remain distinct from the final 5+5 portfolio. Paper~3 has not shown CS predictive power, and Paper~4 has not shown planner superiority. The contribution is a transparent, falsifiable research portfolio with explicit counterarguments and verification criteria. Independent methodological novelty and publication value remain open to external review.

\section*{Author contributions and AI disclosure}
Ada Mao directed and supervised the research program. GPT6-luna, an AI research agent, performed literature synthesis, candidate screening, evidence organization, and prospective qualitative forecast development under human supervision. GPT6-luna is an AI research collaborator, not a human researcher. This manuscript and portfolio remain uncalibrated and have not received independent human peer review.

\section*{Data and reproducibility}
The portfolio index, canonical dossier, hash manifest, and frozen registry are publicly inspectable at the \href{https://github.com/nonameisready/scientific-discovery-computer-science/tree/db6aaa4b7f5fa364ced3042bf55d8df6b8eef1ce}{research repository checkpoint}. The manifest hashes repository artifact bytes and unchanged frozen inputs, not the cited full-text papers. Repository validation checks structure and integrity, not forecast truth. No experiments, model training/inference, new labels, or new forecast commitments were generated in the final portfolio round.

\bibliographystyle{plainnat}
\begin{thebibliography}{99}
\bibitem[Ajith et~al.(2026)]{ajith2026prescience} Ajith, A., et~al. PreScience: A dataset and benchmark for scientific forecasting. arXiv:2602.20459, 2026.
\bibitem[Assadi et~al.(2025)]{assadi2025ors} Assadi, S., Khanna, S., and Kiss, P. Improved bounds for fully dynamic matching via ordered Ruzsa--Szemer\'edi graphs. SODA, 2025. doi:10.1137/1.9781611978322.96.
\bibitem[Behnezhad and Ghafari(2024)]{behnezhad2024dynamic} Behnezhad, S. and Ghafari, M. Fully dynamic matching and ordered Ruzsa--Szemer\'edi graphs. FOCS, 2024. doi:10.1109/FOCS61266.2024.00027.
\bibitem[Bentert et~al.(2024)]{bentert2024twosets} Bentert, M., et~al. Two-Sets Cut-Uncut on planar graphs. ICALP, 2024. doi:10.4230/LIPIcs.ICALP.2024.22.
\bibitem[Bernstein et~al.(2025)]{bernstein2025maximal} Bernstein, A., et~al. Deterministic dynamic maximal matching in sublinear update time. arXiv:2504.20780, 2025.
\bibitem[Brakensiek et~al.(2026)]{brakensiek2026maxbisection} Brakensiek, J., Huang, M., Potechin, A., and Zwick, U. MAX BISECTION might be harder to approximate than MAX CUT. SODA, 2026. doi:10.1137/1.9781611978971.197.
\bibitem[Bravyi et~al.(2024)]{bravyi2024identity} Bravyi, S., Parham, T., and Tran, M.~C. Identity check problem for shallow quantum circuits. arXiv:2401.16525, 2024.
\bibitem[Bringmann et~al.(2019)]{bringmann2019frechet} Bringmann, K., K\"unnemann, M., and Nusser, A. Fr\'echet distance under translation: conditional hardness and an algorithm via offline dynamic grid reachability. arXiv:1810.10982, 2019.
\bibitem[Chen et~al.(2025)]{chen2025subsetsum} Chen, L., Jin, C., Randolph, W., and Servedio, R. New algorithms for pigeonhole equal subset sum. ESA, 2025. doi:10.4230/LIPIcs.ESA.2025.86.
\bibitem[Cslovjecsek et~al.(2025)]{cslovjecsek2025blockip} Cslovjecsek, J., Kouteck\'y, M., Lassota, A., Pilipczuk, M., and Polak, A. Parameterized algorithms for block-structured integer programs with large entries. TheoretiCS 4, 2025. doi:10.46298/theoretics.25.15.
\bibitem[Jaillet et~al.(2025)]{jaillet2025kv} Jaillet, P., et~al. Online scheduling for LLM inference with KV cache constraints. arXiv:2502.07115, 2025.
\bibitem[Jung et~al.(2026)]{jung2026spexis} Jung, H., et~al. Spexis: speculative lookahead scheduling for LLM inference. arXiv:2609.34370, 2026.
\bibitem[Lenzen and Wenning(2024)]{lenzen2024domination} Lenzen, C. and Wenning, S. Tight bounds for constant-round domination on graphs of high girth and low expansion. arXiv:2408.12998, 2024.
\bibitem[Liu et~al.(2026)]{liu2026researchbench} Liu, Y., et~al. ResearchBench: benchmarking LLMs in scientific discovery via inspiration-based task decomposition. Findings of ACL, 2026. doi:10.18653/v1/2026.findings-acl.644.
\bibitem[Mao(2026a)]{mao2026questiondiscovery} Mao, H. Evidence-Based Scientific Question Discovery: A Framework with Historical Backtesting. arXiv:2608.09968, 2026.
\bibitem[Mao(2026b)]{mao2026backtesting} Mao, H. Historical Backtesting for Scientific Question Discovery: A Protocol and Astronomy Pilot. arXiv:2608.16795, 2026.
\bibitem[Min et~al.(2021)]{min2021breakthrough} Min, C., Bu, Y., and Sun, J. Predicting scientific breakthroughs based on knowledge structure variations. Technological Forecasting and Social Change 164, 120502, 2021. doi:10.1016/j.techfore.2020.120502.
\bibitem[Starace et~al.(2025)]{starace2025paperbench} Starace, G., et~al. PaperBench: Evaluating AI's ability to replicate AI research. arXiv:2504.01848, 2025.
\bibitem[Tran et~al.(2026)]{tran2026fdp} Tran, H., et~al. $f$-Differential privacy filters: validity and approximate solutions. arXiv:2602.06756, 2026.
\bibitem[Venkataramani et~al.(2026)]{venkataramani2026masprove} Venkataramani, V., et~al. MAS-ProVe: understanding the process verification of multi-agent systems. arXiv:2602.03053, 2026.
\bibitem[Xin et~al.(2026)]{xin2026verisoft} Xin, Y., et~al. VeriSoftBench: repository-scale formal verification benchmarks for Lean. arXiv:2602.18307, 2026.
\bibitem[Yao et~al.(2026)]{yao2026cocoprover} Yao, S., et~al. Cost-efficient theorem proving via agent orchestration in program verification. arXiv:2610.09681, 2026.
\bibitem[Ye et~al.(2026)]{ye2026vero} Ye, Z., et~al. Vero: can AI agents build formally verified software repositories? arXiv:2608.13522, 2026.
\end{thebibliography}
\end{document}
